\begin{definition}[Grupo]
  Um Grupo $(G,*)$ é um conjunto $G$ não vazio com uma operação $*$ que satisfaz as seguintes propriedades
\begin{itemize}
  \item \textbf{Fechamento:} Se $a,b\in G$, então $a*b\in G$.
  \item \textbf{Associatividade:} Para quaisquer $a,b,c\in G$, tem-se que $a*(b*c)=(a*b)*c$.
  \item \textbf{Existência do Elemento Identidade:} Existe um elemento $e\in G$ tal que $e*a=a=a*e$ para tudo $a\in G$.
  \item \textbf{Existência do Elemento Inverso:} Para cada $a\in G$, existe um elemento $b\in G$ tal que $a*b=e=b*a$. Geralmente o elemento identidade se denota por $e$ o $1$. 
\end{itemize}
\end{definition}
\begin{definition}[Grupo Abeliano]
  Um Grupo $(G,*)$ é um grupo abeliano si ele é Grupo y alem do mais, ele satisfaz
  \begin{itemize}
    \item \textbf{Comutatividade:} Para todo par $a,b\in G$, tem-se que $a*b=b*a$.
  \end{itemize}
\end{definition}
A partir de agora, usaremos notação multiplicativa.
\begin{definition}[Ordem do grupo]
  Definimos ordem do grupo a o número de elementos do conjunto $G$, denotado por $|G|$. 
\end{definition}
\begin{definition}[Ordem do elemento]
  A ordem de um elemento $a\in G$ é definida como o menor enteiro positivo $n$ tal que $a^{n}=\underbrace{a\cdot a\cdots a}_{n \text{ vezes}}=1$ e denotada por $o(a)=n$. Se não existe tal $n$, dizemos que $a$ tem ordem infinita e denotamos $o(a)=\infty$. 
\end{definition}
\begin{example}
  Se pode verificar que $(\mathbb{Z},+)$ é um grupo abeliano. 
\end{example}
\begin{example}
  Se pode verificar que $(\mathbb{R}^{*},\cdot)$ é um grupo abeliano. 
\end{example}
\begin{example}[$GL_n$]
  O conjunto das matrizes invertíveis sobre um corpo $K$ é denotado por $GL_n(K)$ é denominado o \textbf{grupo linear geral de grau $n$ sobre $K$}. Se você suponhe que $n\geq 2$, então
  \begin{align*}
    GL_n(K)=\{A\in M_{n\times n}: \det(A)\neq 0\}
  \end{align*}
  é um grupo não abeliano.
\end{example}
\begin{example}
  Seja $X$ um conjunto não vazio, o conjunto
  \begin{align*}
    Sim(X)=\{f:X\to X : f\text{ e bijeção}\}
  \end{align*}
  é um grupo com a operação de composição de funções. 
\end{example}
\begin{note}[$S_n$]
  Se $X=\{1,\cdots,n\}$ denotamos $Sim(X)$ por $S_n$ (grupo de permutações de $n$ símbolos), dito \textbf{ grupo simétrico de grau $n$}. 
\end{note}
\begin{example}[$D_n$]
  O grupo de simetrias de um polígono regular de $n$ lados é o \textbf{grupo diedral de grau $n$}, denotado por $D_n$, que possui $n$ reflexões y $n$ rotações. $D_n$ é um grupo no abeliano. 
\end{example}
\begin{example}
  Sejan $(G,\star)$ e $(H,\diamond)$ grupos. Defina uma operação em $G\times H$ por
  \begin{align*}
    (g,h)*(g^{\prime},h^{\prime})=(g\star g^{\prime},h\diamond h^{\prime}).
  \end{align*}
  Então $G\times H$ é um grupo. Se $G$ e $H$ são abelianos, então $G\times H$ também o é. Se $G$ e $H$ são finitos, então $G\times H$ também lo é e $|G\times H|=|G||H|$. 
\end{example}
